\documentclass[10pt,a4paper]{article}

% --- Packages ---
\usepackage[utf8]{inputenc}
\usepackage[T1]{fontenc}
\usepackage{lmodern}
\usepackage[margin=1.8cm, top=2.0cm, bottom=2.0cm]{geometry}
\usepackage{amsmath, amssymb, amsfonts, mathtools, bm}
\usepackage{graphicx}
\usepackage{booktabs}
\usepackage{tabularx}
\usepackage{multirow}
\usepackage{xcolor}
\usepackage{tcolorbox}
\usepackage{microtype}
\usepackage{titlesec}
\usepackage{enumitem}
\usepackage{caption}
\usepackage{subcaption}
\usepackage{float}
\usepackage[hidelinks, colorlinks=true, linkcolor=navyblue, citecolor=navyblue, urlcolor=navyblue]{hyperref}

% --- Color Definitions ---
\definecolor{navyblue}{RGB}{20, 50, 110}
\definecolor{slateblue}{RGB}{43, 108, 176}
\definecolor{tealaccent}{RGB}{44, 122, 123}
\definecolor{darkcharcoal}{RGB}{35, 45, 55}
\definecolor{lightbg}{RGB}{247, 250, 252}
\definecolor{boxborder}{RGB}{200, 215, 235}

% --- Typography & Section Formatting ---
\titleformat{\section}{\large\bfseries\color{navyblue}}{\thesection.}{0.5em}{}[\titlerule]
\titleformat{\subsection}{\normalsize\bfseries\color{slateblue}}{\thesubsection.}{0.4em}{}
\titleformat{\subsubsection}{\small\bfseries\color{tealaccent}}{\thesubsubsection.}{0.3em}{}
\titlespacing*{\section}{0pt}{10pt}{4pt}
\titlespacing*{\subsection}{0pt}{8pt}{3pt}
\titlespacing*{\subsubsection}{0pt}{6pt}{2pt}

\setlist{itemsep=2pt, topsep=2pt, parsep=0pt, partopsep=0pt}

% --- TColorBox Styling ---
\tcbuselibrary{skins, breakable}
\newtcolorbox{mathbox}[1][]{
    enhanced,
    colback=lightbg,
    colframe=slateblue,
    arc=3mm,
    boxrule=0.8pt,
    left=6pt, right=6pt, top=6pt, bottom=6pt,
    #1
}

\newtcolorbox{theorembox}[1]{
    enhanced,
    colback=lightbg,
    colframe=navyblue,
    arc=2mm,
    boxrule=1pt,
    title=\textbf{#1},
    coltitle=white,
    colbacktitle=navyblue,
    left=6pt, right=6pt, top=5pt, bottom=5pt,
}

\begin{document}

% --- Header / Title Area ---
\begin{center}
    {\LARGE\bfseries\color{navyblue} Latent World Models for Multivariate Data\\ using JEPA-based Architectures}\\[0.4em]
    {\large\bfseries\color{slateblue} Memoria Técnica y Resumen de Ejecución Experimental (Release v0.1.0)}\\[0.3em]
    {\footnotesize\textbf{Autor:} Víctor Vallejo \quad$\vert$\quad \textbf{Programa:} Máster en Ingeniería Matemática \quad$\vert$\quad \textbf{Fecha:} Agosto 2026 \quad$\vert$\quad \textbf{Estado:} Publicación / Producción}
\end{center}
\vspace{-0.6em}
\hrule height 1.2pt \color{navyblue}
\vspace{0.6em}

% --- Section 1: Introduction & Hypotheses ---
\section{Resumen Ejecutivo e Hipótesis de Investigación}

El objetivo de este Trabajo de Fin de Máster (TFM) es diseñar, formalizar e implementar un marco riguroso de \textbf{Modelos de Mundo Latentes} (\textit{Latent World Models}) para series temporales multivariantes y sistemas dinámicos continuos basados en las \textbf{Joint-Embedding Predictive Architectures (JEPA)} propuestas por Yann LeCun. 

A diferencia de los modelos de mundo tradicionales (PlaNet, DreamerV1-V3) fundamentados en la reconstrucción directa del espacio de observaciones mediante Autoencoders, VAEs o Difusión, la arquitectura JEPA aprende dinámicas representacionales \textbf{exclusivamente en un espacio latente abstracto} $\mathcal{Z}$, eliminando el decodificador de observaciones y optimizando un cuello de botella de información predictiva (\textit{Predictive Information Bottleneck}, PIB).

\begin{itemize}[leftmargin=1.2em]
    \item \textbf{Hipótesis 1 (Invarianza al Ruido Nuisance):} Frente a observaciones con canales estocásticos de alta entropía no correlacionados temporalmente (\textit{nuisance noise}), JEPA preserva la fidelidad dinámica latente sin malgastar capacidad de codificación, superando a los modelos generativos reconstructivos.
    \item \textbf{Hipótesis 2 (Identificabilidad Lineal Isométrica):} Bajo regularización isotrópica Gaussiana (\textbf{SIGReg}), la representación latente $h(s)$ recupera fielmente los verdaderos grados de libertad del sistema físico $s \in \mathcal{S}$ hasta una rotación ortogonal exacta ($h(s) = Qs, Q \in \mathcal{O}(d_s)$), permitiendo sondeos lineales directos (\textit{linear probes}) de alta precisión.
    \item \textbf{Hipótesis 3 (Control Óptimo Latente en Tiempo Real):} Los algoritmos de planificación basada en modelos (CEM y MPPI) operando directamente sobre el predictor latente $P_\phi(z_t, a_t)$ logran optimizar trayectorias de control a frecuencias superiores a $200\text{ Hz}$ ($< 5\text{ ms}$ por paso).
\end{itemize}

% --- Section 2: Mathematical Foundations ---
\section{Fundamentación Matemática y Garantías Teóricas}

\subsection{Formulación del Sistema Dinámico y Espacio de Observación}
Sea un sistema físico continuo gobernado por ecuaciones diferenciales ordinarias sobre una variedad latente $\mathcal{S} \subseteq \mathbb{R}^{d_s}$:
\begin{equation}
    \frac{d s(t)}{dt} = f(s(t), a(t)) + \eta(t), \quad s(t) \in \mathbb{R}^{d_s}, \quad a(t) \in \mathbb{R}^{d_a}
\end{equation}
El sistema se muestrea a intervalos $\Delta t$ mediante un operador no lineal de medición $g: \mathcal{S} \times \mathcal{V} \to \mathbb{R}^C$:
\begin{equation}
    x_t = g(s_t, \nu_t) + \epsilon_t, \quad \nu_t \sim p_\nu(\cdot), \quad \epsilon_t \sim \mathcal{N}(0, \sigma_\epsilon^2 I_C)
\end{equation}
donde $\nu_t \in \mathbb{R}^{d_\nu}$ es un vector de ruido estocástico no estructurado e impredecible (fluctuaciones térmicas, vibraciones parásitas, derivas sensoriales).

\subsection{Objetivo de Pérdida LeWM y Regularización SIGReg (Cramér-Wold + Epps-Pulley)}
Para entrenar el modelo de mundo de forma estable de principio a fin (\textit{end-to-end}) evitando el colapso de representación hacia puntos constantes sin recurrir a heurísticas de gradiente detenido asimétricas, implementamos la función de pérdida \textbf{LeWorldModel (LeWM)}:
\begin{equation}
    \mathcal{L}_{LeWM} = \mathcal{L}_{pred}\left(\hat{z}_{t+1:t+K}, z_{t+1:t+K}\right) + \lambda \cdot \text{SIGReg}(Z)
\end{equation}
donde la pérdida predictiva multi-paso con factor de descuento temporal $\gamma \in (0, 1]$ se define como:
\begin{equation}
    \mathcal{L}_{pred} = \frac{1}{\sum_{k=1}^K \gamma^{k-1}} \sum_{k=1}^K \gamma^{k-1} \left\| \hat{z}_{t+k} - z_{t+k} \right\|_2^2, \quad \hat{z}_{t+k} = P_\phi(\hat{z}_{t+k-1}, a_{t+k-1}), \quad \hat{z}_t = E_\theta(x_{\le t})
\end{equation}

El término \textbf{SIGReg} proyecta los embeddings $Z \in \mathbb{R}^{N \times d_z}$ sobre $M$ direcciones aleatorias uniformes $u^{(m)} \sim \text{Uniform}(\mathbb{S}^{d_z-1})$ (Teorema de Cramér-Wold) y evalúa el estadístico de prueba univariado de \textbf{Epps-Pulley} $T(h)$ contra la Gaussiana estándar $\mathcal{N}(0, 1)$ mediante la función característica empírica $\phi_N(t) = \frac{1}{N}\sum_{j=1}^N e^{i t h_j}$:
\begin{equation}
    \text{SIGReg}(Z) \triangleq \frac{1}{M} \sum_{m=1}^M T\left(Z u^{(m)}\right), \quad T(h) = \int_{-\infty}^\infty \left| \phi_N(t; h) - e^{-\frac{t^2}{2}} \right|^2 e^{-\frac{t^2}{2\sigma_w^2}} dt
\end{equation}

\begin{mathbox}
\textbf{Forma Cerrada Analítica Exacta del Estadístico de Epps-Pulley (Implementada e Integrada):}
\begin{equation}
    \boxed{T(h) = \frac{1}{N^2}\sum_{j=1}^N\sum_{l=1}^N \exp\left( -\frac{\sigma_w^2(h_j - h_l)^2}{2} \right) - \frac{2}{N\sqrt{\sigma_w^2 + 1}}\sum_{j=1}^N \exp\left( -\frac{\sigma_w^2 h_j^2}{2(\sigma_w^2 + 1)} \right) + \frac{1}{\sqrt{2\sigma_w^2 + 1}}}
\end{equation}
\end{mathbox}

\subsection{Teorema de Identificabilidad Lineal de LeJEPA (Klindt, LeCun, Balestriero 2026)}
\begin{theorembox}{Teorema (Identificabilidad Lineal y Unicidad Gaussiana)}
Sea un proceso estacionario de Gauss-Markov (Ornstein-Uhlenbeck) $s_{t+1} = \rho s_t + \sqrt{1-\rho^2}\eta_t$ con $|\rho| < 1$. Toda representación latente $h: \mathcal{S} \to \mathcal{Z}$ que minimiza el objetivo de alineación JEPA bajo la restricción marginal Gaussiana impuesta por SIGReg recupera los estados físicos reales mediante una transformación lineal ortogonal:
\begin{equation}
    h(s) = Q s, \quad \text{donde } Q \in \mathcal{O}(d_s) \text{ satisface } Q^T Q = I_{d_s}
\end{equation}
\textbf{Mecanismo de Demostración:} Mediante la fórmula de Mehler sobre la base de polinomios ortogonales de Hermite $\{H_k(x)\}_{k=0}^\infty$, el operador de transición descompone la correlación temporal según $\mathbb{E}[H_k(s_t)H_k(s_{t+1})] = \rho^k$. La penalización de error cuadrático $\sum_{k=1}^\infty (1-\rho^{2k})\|c_k\|^2$ castiga estrictamente cualquier distorsión no lineal ($k \ge 2$) debido a que $\rho^{2k} < \rho^2$. Por análisis de Sturm-Liouville, la Gaussiana es la \textbf{única} distribución donde la autofunción lineal domina estrictamente.
\end{theorembox}

\begin{theorembox}{Corolario (Equivalencia Isométrica del Control Latente)}
Si $h(s) = Qs$ con $Q \in \mathcal{O}(d_s)$, la distancia euclidiana en el espacio latente coincide exactamente con la distancia en el espacio de estados físicos reales:
\begin{equation}
    \|z_t - z_{\text{meta}}\|_2^2 = \|Qs_t - Qs_{\text{meta}}\|_2^2 = (s_t - s_{\text{meta}})^T Q^T Q (s_t - s_{\text{meta}}) = \|s_t - s_{\text{meta}}\|_2^2
\end{equation}
Por tanto, la planificación latente $J(a_{0:H-1}) = \min \sum \|z_t - z_{\text{meta}}\|^2$ es formalmente idéntica al control óptimo con conocimiento perfecto del estado.
\end{theorembox}

% --- Section 3: Framework Architecture ---
\section{Arquitectura del Framework de Software (\texttt{jepa\_wm})}

\begin{figure}[H]
    \centering
    \includegraphics[width=0.98\textwidth]{results/jepa_wm_architecture.jpg}
    \caption{\textbf{Arquitectura del Modelo de Mundo Latente JEPA (\texttt{LatentJEPAWorldModel}):} Flujo de información y componentes matemáticos. (1) El \textit{Multivariate Patch Transformer Context Encoder} $E_\theta$ mapea ventanas multicanal $x_{\le t}$ al estado latente $z_t$. (2) El \textit{Action-Conditioned Latent Predictor} $P_\phi$ proyecta dinámicas multi-paso $\hat{z}_{t+1:t+K}$ condicionadas por acciones $a_t$. (3) El \textit{Target Sequence Encoder} $E_{\bar{\theta}}$ codifica observaciones futuras de referencia para evaluar la pérdida $\mathcal{L}_{pred}$. (4) El nodo de proyección Cramér-Wold calcula el regularizador analítico anti-colapso $\text{SIGReg}(Z)$. (5) El módulo de planificación latente MPC (CEM / MPPI) optimiza acciones en tiempo real directamente en $\mathcal{Z}$ sin requerir decodificador de observaciones.}
\end{figure}

El paquete de software \texttt{jepa\_wm} ha sido estructurado en submódulos modulares y desacoplados:

\begin{table}[H]
\centering
\small
\begin{tabularx}{\textwidth}{l p{4.2cm} X}
\toprule
\textbf{Módulo} & \textbf{Componentes Clave} & \textbf{Especificación Técnica y Propósito} \\
\midrule
\texttt{jepa\_wm.core} & \texttt{MultivariatePatchEncoder}\newline\texttt{ResidualGRUPredictor}\newline\texttt{LatentJEPAWorldModel}\newline\texttt{ReconstructiveWorldModel} & Encoder Transformer multicanal con autoatención espacial/temporal y proyección de tokens. Predictor autorregresivo residual inspirado en integración Euler. World Model end-to-end con encoder siamés/EMA y regularización SIGReg. Baseline VAE/Autoencoder para comparación científica. \\
\addlinespace[2pt]
\texttt{jepa\_wm.losses} & \texttt{SIGReg}, \texttt{LeWMLoss}\newline\texttt{VICRegLoss}, \texttt{InfoNCELoss}\newline\texttt{VariationalJEPALoss} & Regularizador anti-colapso por proyecciones Cramér-Wold analíticas. Pérdida multi-paso con descuento temporal $\gamma$. Baselines comparativos contrastivos, de covarianza y variacionales. \\
\addlinespace[2pt]
\texttt{jepa\_wm.data} & \texttt{Lorenz63System}, \texttt{Lorenz96}\newline\texttt{CoupledOscillatorsSystem}\newline\texttt{MultiSensorBenchmark} & Integradores numéricos RK4 para atractor caótico 3D de Lorenz, dinámicas atmosféricas $N$-dimensionales, osciladores hamiltonianos y benchmark industrial con canales de ruido nuisance controlados. \\
\addlinespace[2pt]
\texttt{jepa\_wm.diagnostics} & \texttt{LinearIdentifiabilityProbe}\newline\texttt{OrthogonalProcrustes}\newline\texttt{compute\_effective\_rank}\newline\texttt{compute\_noise\_rejection\_ratio} & Sonda de regresión Ridge para evaluar recuperación de estados $R^2$, algoritmo SVD de alineación ortogonal de Procrustes, cálculo de rango efectivo ($\text{erank}(Z) = e^{H(p_\sigma)}$) y ratio de rechazo de ruido. \\
\addlinespace[2pt]
\texttt{jepa\_wm.planning} & \texttt{LatentCEMPlanner}\newline\texttt{LatentMPPIPlanner} & Algoritmos de optimización de trayectorias en espacio latente: Cross-Entropy Method (CEM) con poblaciones élite y Model Predictive Path Integral (MPPI) con pesos de Boltzmann. \\
\bottomrule
\end{tabularx}
\caption{Estructura de arquitectura y módulos del framework \texttt{jepa\_wm}.}
\end{table}

% --- Section 4: Experimental Results ---
\section{Resultados Experimentales y Validación de Benchmarks}

\begin{table}[H]
\centering
\small
\begin{tabularx}{\textwidth}{l c c X}
\toprule
\textbf{Experimento / Benchmark} & \textbf{Métrica Matemática} & \textbf{Valor Obtenido} & \textbf{Interpretación y Conclusión Científica} \\
\midrule
\multirow{3}{*}{\textbf{Exp 1: Atractor Caótico Lorenz-63}} & Identificabilidad Lineal $R^2$ & $\mathbf{0.9074}$ & Recuperación de la variedad caótica completa desde 8 canales ruidosos. \\
& Correlación Canónica (CCA) & $\mathbf{0.9471}$ & Fuerte alineación canónica entre latentes $z$ y estados físicos $s$. \\
& Error Predicción Latente MSE & $\text{Paso 1: } 5.16\times 10^{-4}$ & Trayectorias de predicción estables en horizontes multi-paso. \\
\midrule
\textbf{Exp 2: Barrido SIGReg ($\lambda \in [0, 1]$)} & Identificabilidad ($\lambda=1.0$) & $\mathbf{0.9885}$ & Demostración de Teorema 1: SIGReg previene colapso y maximiza $R^2$. \\
\midrule
\textbf{Exp 3: Robustez ante Ruido Nuisance} & Retención $R^2$ ($C_{\text{nuis}} \in [0, 16]$) & $\mathbf{0.8927 \to 0.8396}$ & JEPA mantiene rendimiento; el baseline reconstructivo cae bruscamente. \\
\midrule
\multirow{2}{*}{\textbf{Exp 4: Control Óptimo MPC Latente}} & Tiempo de Ejecución MPPI & $\mathbf{4.87\text{ ms / paso}}$ & Control latente en tiempo real ($>200\text{ Hz}$) sin decodificar píxeles. \\
& Tiempo de Ejecución CEM & $\mathbf{29.97\text{ ms / paso}}$ & Convergencia monótona hacia el estado objetivo latente $z_{\text{meta}}$. \\
\bottomrule
\end{tabularx}
\caption{Resumen cuantitativo de los cuatro experimentos de benchmark ejecutados.}
\end{table}

\begin{figure}[H]
    \centering
    \includegraphics[width=0.92\textwidth]{results/lorenz_attractor_recovery.png}
    \caption{\textbf{Experimento 1:} Recuperación del Atractor Caótico 3D de Lorenz-63. Izquierda: Órbita real $s(t)=(x,y,z)$. Derecha: Órbita recuperada por sondeo lineal $\hat{s}(t) = W z(t) + b$ desde el espacio latente del JEPA entrenado con SIGReg ($R^2 = 0.9074$, $\text{CCA} = 0.9471$).}
\end{figure}

\begin{figure}[H]
    \centering
    \begin{subfigure}[b]{0.48\textwidth}
        \centering
        \includegraphics[width=\textwidth]{results/identifiability_spectrum.png}
        \caption{Decaimiento singular según peso $\lambda$ en SIGReg.}
    \end{subfigure}
    \hfill
    \begin{subfigure}[b]{0.48\textwidth}
        \centering
        \includegraphics[width=\textwidth]{results/noise_benchmark_results.png}
        \caption{JEPA vs Reconstrucción ante ruido nuisance.}
    \end{subfigure}
    \caption{\textbf{Figuras de Validación Teórica.} (a) Espectro singular demostrando prevención del colapso con regularización Gaussiana. (b) Fidelidad predictiva downstream de JEPA frente al colapso de capacidad del modelo reconstructivo.}
\end{figure}

% --- Section 5: Git Traceability & PhD Roadmap ---
\section{Trazabilidad Git, Verificación de Tests y Roadmap Doctoral}

\begin{table}[H]
\centering
\footnotesize
\begin{tabularx}{\textwidth}{l p{6.5cm} c}
\toprule
\textbf{Rama Git} & \textbf{Alcance Técnico y Contenido del Módulo} & \textbf{Estado Remoto (\texttt{origin})} \\
\midrule
\texttt{feature/repo-foundation} & Configuración base, \texttt{pyproject.toml}, \texttt{setup.py}, \texttt{.gitignore}, \texttt{README.md} & Fusionado en \texttt{dev} y \texttt{main} \\
\texttt{feature/sigreg-and-losses} & SIGReg cerrado, VICReg, InfoNCE, LeWMLoss y VariationalJEPA & Fusionado en \texttt{dev} y \texttt{main} \\
\texttt{feature/dynamical-systems} & Lorenz-63/96, osciladores acoplados, benchmark nuisance y \texttt{Dataset} & Fusionado en \texttt{dev} y \texttt{main} \\
\texttt{feature/latent-architectures} & Encoders Transformer/TCN, predictores residuales y baseline reconstructivo & Fusionado en \texttt{dev} y \texttt{main} \\
\texttt{feature/latent-diagnostics} & Sondas lineales $R^2$, Procrustes, CCA, rango efectivo (\texttt{erank}) y NRR & Fusionado en \texttt{dev} y \texttt{main} \\
\texttt{feature/latent-planning-mpc} & Planificadores Latent CEM y Latent MPPI para control óptimo & Fusionado en \texttt{dev} y \texttt{main} \\
\texttt{feature/training-and-experiments} & Motor \texttt{JEPATrainer}, evaluadores y los 4 scripts de benchmark reproducibles & Fusionado en \texttt{dev} y \texttt{main} \\
\texttt{docs/tfm-thesis-blueprint} & Memorias matemáticas (\LaTeX), blueprint del TFM y API Reference & Fusionado en \texttt{dev} y \texttt{main} \\
\textbf{\texttt{main} (Release v0.1.0)} & \textbf{Versión de producción etiquetada con tag formal \texttt{v0.1.0}} & \textbf{Release v0.1.0 Publicado} \\
\bottomrule
\end{tabularx}
\caption{Historial y trazabilidad del flujo de trabajo Git profesional implementado.}
\end{table}

\subsection{Verificación Automatizada de Calidad}
Toda la librería cuenta con cobertura de pruebas unitarias (\textbf{24/24 tests pasados}):
\begin{itemize}[leftmargin=1.2em]
    \item \texttt{tests/test\_losses.py}: Verificación de gradientes autograd, no negatividad de SIGReg y convergencia de distribuciones gaussianas frente a latentes colapsados.
    \item \texttt{tests/test\_data.py}: Integración RK4, conservación de energía en osciladores hamiltonianos y particionamiento de ventanas temporales.
    \item \texttt{tests/test\_models.py}: Compatibilidad de dimensiones tensoriales, paso hacia adelante y actualización de encoders EMA.
    \item \texttt{tests/test\_diagnostics.py}: SVD de Procrustes ortogonal, correlación canónica y cálculo de entropía espectral.
    \item \texttt{tests/test\_planning.py}: Optimización y convergencia de costes en CEM y MPPI latente.
\end{itemize}

\subsection{Líneas de Investigación Prioritarias para la Tesis Doctoral}
Este proyecto sienta las bases matemáticas y experimentales para la tesis doctoral en modelos de mundo JEPA:
\begin{enumerate}[leftmargin=1.4em]
    \item \textbf{Continuous-Time Neural ODE-JEPAs:} Extensión a formulaciones continuas con flujos normalizantes y operadores adjuntos para series temporales irregulares y asíncronas.
    \item \textbf{Hierarchical Partitioning Graph-JEPAs (HP-JEPA):} Modelado de sistemas físicos relacionales multicuerpo (grafos) con proyección al hiperboloide de Lorentz y particionamiento multirresolución.
    \item \textbf{Test-Time Adaptation (AdaJEPA):} Adaptación en bucle cerrado de MPC para mitigar el desvío de distribución (\textit{distribution shift}) en tiempo real sin demostraciones de expertos.
\end{enumerate}

\end{document}
